\subsection{Removing the degree factor from the temperature condition}
\label{resolution-dynamics:weighted-maps}

Threshold maps give a simple proof, but they may inspect every neighbor
whenever they are nonconstant. A different exact update representation
controls the expected number of inspected neighbors by the incident
interaction strength itself.

\begin{theorem}[A degree-independent sufficient temperature]
\label{thm:resolution-dynamics-weighted}
Fix \(D<\infty\) and \(k\ge2\). Let
\[
 H(x)=-\sum_A J_A\prod_{i\in A}x_i
\]
have interactions of size at most \(k\), arbitrary one-spin fields,
and an interaction graph of maximum degree at most \(D\), where two
spins are adjacent when they share a nonzero interaction. Set
\[
 K=\beta\max_i\sum_{A\ni i,\,|A|\ge2}|J_A|.
\]
If \(K\le1/[8(k-1)]\), the Gibbs laws have hereditary full AA with
query bound \(\operatorname{poly}_{D,k}(T,1/e)\), uniformly over all
finite fields and all model parameters satisfying this inequality.
For pair interactions, the sufficient condition is simply
\(\beta\Lambda\le1/8\).
\end{theorem}

The temperature bound does not depend on graph degree, but the query
polynomial still may: this theorem does not cover a degree bound growing
with input dimension.

\textbf{Two elementary norms.}
For a function on a finite Boolean cube, write
\(f=\sum_S\widehat f(S)\chi_S\), and set
\[
 \|f\|_A=\sum_S|\widehat f(S)|,\qquad
 W(f)=\sum_S|S|\,|\widehat f(S)|.
\]
Since \(\chi_S\chi_T=\chi_{S\triangle T}\),
\[
 \|fg\|_A\le\|f\|_A\|g\|_A,\qquad
 W(fg)\le W(f)\|g\|_A+\|f\|_AW(g).
 \tag{WA}
\]
All Fourier coefficients and norms here concern only the at most
\(D\) neighbors of one site.

\begin{lemma}[An exact sparse-support update law]
\label{lem:resolution-dynamics-parity-law}
Let the conditional mean of one spin be
\(m(x)=\tanh(h+u(x))\), where
\(u\) has no constant term,
\(\|u\|_A\le K\le1/8\), and
\(W(u)\le(k-1)K\). There is an exact random map with conditional
mean \(m(x)\) that is either constant or a signed parity of a subset
of the neighbors. If \(S\) is that subset (empty for a constant map),
then, conditional on any neighbor assignment and either resulting spin,
\[
 \mathbb E[|S|\mid\text{neighbors, resulting spin}]
 \le (k-1)\chi(K),\qquad
 \chi(K)=\frac{4e^{2K}K\sec^2K}{(1-\tan K)^2}.
 \tag{PS}
\]
\end{lemma}

\begin{proof}
Write \(r=\tanh h\) and \(v=\tanh u\). The absolute values of
the Taylor coefficients of \(\tanh z\) are the positive Taylor
coefficients of \(\tan z\). This follows from
\(\tanh z=-i\tan(iz)\); positivity of the latter coefficients also
follows recursively from \((\tan z)'=1+(\tan z)^2\). The series
converge absolutely in the norm \(\|\cdot\|_A\) for \(K<\pi/2\).
Using (WA) termwise gives
\[
 \|v\|_A\le t:=\tan K,
 \qquad W(v)\le(k-1)K\sec^2 K.
\]
The addition formula and a convergent geometric series yield
\[
 m=r+(1-r^2)\frac{v}{1+rv},\qquad
 \|m-r\|_A\le\frac{(1-r^2)t}{1-|r|t},\qquad
 W(m)\le\frac{(1-r^2)(k-1)K\sec^2K}{(1-|r|t)^2}.
 \tag{PM}
\]

Write \(m=\sum_S c_S\chi_S\) and
\(A=\sum_{S\ne\varnothing}|c_S|\). Since \(t\le1/3\),
\[
 |c_\varnothing|+A
 \le |r|+\frac{(1-r^2)t}{1-|r|t}\le1.
\]
Thus the following probabilities define an actual random map:
choose \(\operatorname{sign}(c_S)\chi_S\) with probability
\(|c_S|\) for each nonempty \(S\), and choose the constant maps
\(+1\) and \(-1\) with respective probabilities
\((1+c_\varnothing-A)/2\) and
\((1-c_\varnothing-A)/2\). Its pointwise mean is exactly \(m\),
so it is the heat-bath update law.

Every possible output spin, on every neighbor assignment, has
probability at least
\[
 \frac{1-\tanh(|h|+K)}2
 \ge\frac{(1-|r|)e^{-2K}}2.
\]
The unconditioned expected support size is \(W(m)\). Conditioning
on any particular output can increase its expectation by at most the
inverse of that output probability. Divide the last bound in (PM) by
this probability and use \(1+|r|\le2\) and
\(1-|r|t\ge1-t\) to obtain (PS).
\end{proof}

\textbf{Uniform contraction constant.}
For \(0\le K\le1/8\), the elementary inequalities
\[
 e^{2K}\le4/3,\qquad \sec^2K\le16/15,
 \qquad\tan K\le2/15
\]
give \(\chi(K)\le(1280/169)K\). For example,
\(\cos K\ge1-K^2/2\ge127/128\) and \(\sin K\le K\)
prove the tangent bound; \(\cos^2K\ge1-K^2\) proves the secant
bound. Consequently the theorem's assumption gives
\[
 \vartheta:=(k-1)\chi(K)\le160/169<1.
 \tag{PC}
\]

\textbf{Proof of Theorem~\ref{thm:resolution-dynamics-weighted}.}
At every site apply the lemma to its interaction field
\(u_i=\beta\sum_{A\ni i,\,|A|\ge2}J_A\chi_{A\setminus\{i\}}\).
It satisfies the stated norm bounds. Use independent parity-map
choices in the same proper-color schedule as before. There are at
most \(2^D-1\) nonconstant maps and two constant maps, again encoded
by \(D+1\) bits of an arbitrary source-block law. Conditional on the
complete spin trajectory these map blocks are independent, just as
for the threshold maps.

To trace a disagreement backwards through an event, only neighbors in
its selected subset \(S\) need be considered. Conditional on the
trajectory, the expected number of such parents is at most
\(\vartheta\), uniformly at every event. Expanding all possible
dependency paths, the expected number of surviving paths of length
\(L\) is at most \(\vartheta^L\): along each path event times
strictly decrease, so its map blocks are distinct and independent
under this conditioning; summing over each next parent uses (PS).
Repeated vertices reached by different paths do not affect this
first-moment argument. Any dependence on the initial configuration
after \(L\) whole sweeps requires such a path. Averaging over
trajectories conditional on the full final input therefore gives
\[
 \Pr[F_L(Z)_i\ne X_i\mid X=x]\le\vartheta^L.
\]

All deterministic forward cones and the exact reverse-history
posterior oracle from the previous proof are unchanged. In particular,
sampling a forward parity map conditional on its local neighbors and
new spin uses only those local values, regardless of how small its
prior weight is. Choose \(L\) using the fixed bound \(160/169\),
and apply the same query bound (QC). It is a polynomial for fixed
\(D,k\). Pinning can only reduce interaction sizes, primal degree,
and the sum of incident nonconstant coupling magnitudes; terms that
become fields are allowed freely. Thus the assertion is uniform over
pinnings as well. \qed
